%% LyX 2.5.1 created this file.  For more info, see https://www.lyx.org/.
%% Do not edit unless you really know what you are doing.
\documentclass[oneside,english]{amsart}
\usepackage[T1]{fontenc}
\usepackage[latin9]{inputenc}
\synctex=-1
\usepackage{xcolor}
\colorlet{note_fontcolor}{lightgray}
\usepackage{enumitem}
\usepackage{amsthm}

\makeatletter

%%%%%%%%%%%%%%%%%%%%%%%%%%%%%% LyX specific LaTeX commands.
%% The greyedout annotation environment
\newenvironment{lyxgreyedout}
{\normalfont\normalsize\textcolor{note_fontcolor}\bgroup\ignorespaces}
{\ignorespacesafterend\egroup}

%%%%%%%%%%%%%%%%%%%%%%%%%%%%%% Textclass specific LaTeX commands.
\newlength{\lyxlabelwidth}      % auxiliary length 

%%%%%%%%%%%%%%%%%%%%%%%%%%%%%% User specified LaTeX commands.
\date{}
\usepackage{commath}

\makeatother

\theoremstyle{plain}
\newtheorem{thm}{\protect\theoremname}
\newtheorem{prop}[thm]{\protect\propositionname}
\theoremstyle{remark}
\newtheorem*{rem*}{\protect\remarkname}
\newtheorem*{claim*}{\protect\claimname}
\usepackage{babel}
\providecommand{\claimname}{Claim}
\providecommand{\propositionname}{Proposition}
\providecommand{\remarkname}{Remark}
\providecommand{\theoremname}{Theorem}

\begin{document}

\section{Counting Polygon Mosaics}

\label{section: counting polygon mosaics}
\begin{thm}
Let $v$ be any cell product function. Define the $1\times1$ matrices
$A_{1}=\begin{pmatrix}v\left(\begin{smallmatrix}0 & 0\\
0 & 0
\end{smallmatrix}\right)\end{pmatrix}$, $B_{1}=\begin{pmatrix}v\left(\begin{smallmatrix}0 & 0\\
1 & 0
\end{smallmatrix}\right)\end{pmatrix}$, $C_{1}=\begin{pmatrix}v\left(\begin{smallmatrix}1 & 0\\
0 & 0
\end{smallmatrix}\right)\end{pmatrix}$, and $D_{1}=\begin{pmatrix}v\left(\begin{smallmatrix}1 & 0\\
1 & 0
\end{smallmatrix}\right)\end{pmatrix}$, and for positive integers $n$, define the $2^{n}\times2^{n}$ matrices
\begin{lyxgreyedout}
Note that in the expressions below like $v(\begin{smallmatrix}0 & 0\\
0 & 1
\end{smallmatrix})B_{n}$, the LaTeX should be updated to \textbackslash left( and \textbackslash right).%
\end{lyxgreyedout}
\begin{eqnarray*}
A_{n+1}=\begin{pmatrix}v(\begin{smallmatrix}0 & 0\\
0 & 0
\end{smallmatrix})A_{n} & v(\begin{smallmatrix}0 & 0\\
0 & 1
\end{smallmatrix})B_{n}\\
v(\begin{smallmatrix}0 & 1\\
0 & 0
\end{smallmatrix})C_{n} & v(\begin{smallmatrix}0 & 1\\
0 & 1
\end{smallmatrix})D_{n}
\end{pmatrix} & B_{n+1}=\begin{pmatrix}v(\begin{smallmatrix}0 & 0\\
1 & 0
\end{smallmatrix})A_{n} & v(\begin{smallmatrix}0 & 0\\
1 & 1
\end{smallmatrix})B_{n}\\
v(\begin{smallmatrix}0 & 1\\
1 & 0
\end{smallmatrix})C_{n} & v(\begin{smallmatrix}0 & 1\\
1 & 1
\end{smallmatrix})D_{n}
\end{pmatrix}\\
C_{n+1}=\begin{pmatrix}v(\begin{smallmatrix}1 & 0\\
0 & 0
\end{smallmatrix})A_{n} & v(\begin{smallmatrix}1 & 0\\
0 & 1
\end{smallmatrix})B_{n}\\
v(\begin{smallmatrix}1 & 1\\
0 & 0
\end{smallmatrix})C_{n} & v(\begin{smallmatrix}1 & 1\\
0 & 1
\end{smallmatrix})D_{n}
\end{pmatrix} & D_{n+1}=\begin{pmatrix}v(\begin{smallmatrix}1 & 0\\
1 & 0
\end{smallmatrix})A_{n} & v(\begin{smallmatrix}1 & 0\\
1 & 1
\end{smallmatrix})B_{n}\\
v(\begin{smallmatrix}1 & 1\\
1 & 0
\end{smallmatrix})C_{n} & v(\begin{smallmatrix}1 & 1\\
1 & 1
\end{smallmatrix})D_{n}
\end{pmatrix}.
\end{eqnarray*}
Then for all positive integers $m$ and $n$, the upper left entry
of $A^{m}_{n}$ equals $\sum_{\ell}v\left(\ell\right)$, where the
sum is over all $\left(m,n\right)$ framed lattices $\ell$. \label{thm: m by n matrix} 
\end{thm}

In the statement of Theorem~\ref{thm: m by n matrix}, we evaluate
$v$ at each of the sixteen possible cells at least once. The order
in which these values appear may look haphazard but is exactly what
we need to establish the first part of Proposition~\ref{prop: one by n matrix},
which we will use to prove the proposition's second part, which in
turn is a key part of the proof of Theorem~\ref{thm: m by n matrix}.

Until we complete the proof of Theorem~\ref{thm: m by n matrix},
we use, as found there, $v$ and the matrices $A_{n},B_{n},C_{n},$
and $D_{n}$, for $n\ge1$. Also, let $\beta_{n}(k)$ be the $n$
digit binary representation of the nonnegative integer $k<2^{n}$,
separated out into individual digits (e.g., $\beta_{3}\left(5\right)$
is $1\ 0\ 1$ rather than $101$), with $\beta_{0}\left(0\right)$
being the empty string. This allows us, for example, to write the
binary lattice $\begin{smallmatrix}0 & 1 & 1 & 0\\
0 & 0 & 1 & 0
\end{smallmatrix}$ as $\begin{smallmatrix}\beta_{4}\left(6\right)\\
\beta_{4}\left(2\right)
\end{smallmatrix}$ or as $\begin{smallmatrix}0 & \beta_{2}\left(3\right) & 0\\
0 & \beta_{2}\left(1\right) & 0
\end{smallmatrix}$.
\begin{prop}
\label{prop: one by n matrix} Let $n$ be a positive integer. Numbering
rows and columns beginning at $0$, we have the following.
\begin{enumerate}
\item $\begin{pmatrix}A_{n} & B_{n}\\
C_{n} & D_{n}
\end{pmatrix}$ equals the $2^{n}\times2^{n}$ matrix whose $\left(i,j\right)$ entry
is $v\left(\begin{smallmatrix}\beta_{n+1}\left(2i\right)\\
\beta_{n+1}\left(2j\right)
\end{smallmatrix}\right)$.
\item $A_{n}$ equals the $2^{n-1}\times2^{n-1}$ matrix whose $\left(i,j\right)$
entry is $v\left(\begin{smallmatrix}0 & \beta_{n-1}\left(i\right) & 0\\
0 & \beta_{n-1}\left(j\right) & 0
\end{smallmatrix}\right)$.
\end{enumerate}
\end{prop}

\begin{proof}
Notice that for all nonnegative integers $k<2^{n}$, $\beta_{n+1}\left(2k\right)$
is the string $\beta_{n}\left(k\right)\ 0$ and that if $k<2^{n-1}$
then $\beta_{n}\left(k\right)$ is the string $0\ \beta_{n-1}\left(k\right)$.
Putting these together, the first part of the proposition implies
the second.

Define $X_{n}$ to be the $2^{n}\times2^{n}$ matrix $\begin{pmatrix}A_{n} & B_{n}\\
C_{n} & D_{n}
\end{pmatrix}$ and $Y_{n}$ to be the $2^{n}\times2^{n}$ matrix whose $\left(i,j\right)$
entry is $v\left(\begin{smallmatrix}\beta_{n+1}\left(2i\right)\\
\beta_{n+1}\left(2j\right)
\end{smallmatrix}\right)$. To complete the proposition, we must prove $X_{n}=Y_{n}$ for all
positive integers $n$. Since $\beta_{2}\left(0\right)$ is $0\ 0$
and $\beta_{2}\left(2\right)$ is $1\ 0$, the $n=1$ case follows
from the definitions of $A_{1}$, $B_{1}$, $C_{1}$, and $D_{1}$
in Theorem~\ref{thm: m by n matrix}. Continuing by induction, we
assume $X_{n}=Y_{n}$ and prove $X_{n+1}=Y_{n+1}$.

The statement of Theorem~\ref{thm: m by n matrix} tells us that
\[
X_{n+1}=\begin{pmatrix}v\left(\begin{smallmatrix}0 & 0\\
0 & 0
\end{smallmatrix}\right)A_{n} & v\left(\begin{smallmatrix}0 & 0\\
0 & 1
\end{smallmatrix}\right)B_{n} & v\left(\begin{smallmatrix}0 & 0\\
1 & 0
\end{smallmatrix}\right)A_{n} & v\left(\begin{smallmatrix}0 & 0\\
1 & 1
\end{smallmatrix}\right)B_{n}\\
v\left(\begin{smallmatrix}0 & 1\\
0 & 0
\end{smallmatrix}\right)C_{n} & v\left(\begin{smallmatrix}0 & 1\\
0 & 1
\end{smallmatrix}\right)D_{n} & v\left(\begin{smallmatrix}0 & 1\\
1 & 0
\end{smallmatrix}\right)C_{n} & v\left(\begin{smallmatrix}0 & 1\\
1 & 1
\end{smallmatrix}\right)D_{n}\\
v\left(\begin{smallmatrix}1 & 0\\
0 & 0
\end{smallmatrix}\right)A_{n} & v\left(\begin{smallmatrix}1 & 0\\
0 & 1
\end{smallmatrix}\right)B_{n} & v\left(\begin{smallmatrix}1 & 0\\
1 & 0
\end{smallmatrix}\right)A_{n} & v\left(\begin{smallmatrix}1 & 0\\
1 & 1
\end{smallmatrix}\right)B_{n}\\
v\left(\begin{smallmatrix}1 & 1\\
0 & 0
\end{smallmatrix}\right)C_{n} & v\left(\begin{smallmatrix}1 & 1\\
0 & 1
\end{smallmatrix}\right)D_{n} & v\left(\begin{smallmatrix}1 & 1\\
1 & 0
\end{smallmatrix}\right)C_{n} & v\left(\begin{smallmatrix}1 & 1\\
1 & 1
\end{smallmatrix}\right)D_{n}
\end{pmatrix},
\]
so that each entry of $X_{n+1}$ is $v$ evaluated at a cell times
an entry of $A_{n}$, $B_{n}$, $C_{n}$, or $D_{n}$. After we also
express every entry of $Y_{n+1}$ as $v$ evaluated at a cell times
another quantity, we will show that the cells evaluated at $v$ are
the same, as are the other quantities, and thus $X_{n+1}=Y_{n+1}$.
To express an entry of $Y_{n+1}$ in our desired way, we use that
$v$ is a cell product function to write $v\left(\begin{smallmatrix}\beta_{n+2}\left(2i\right)\\
\beta_{n+2}\left(2j\right)
\end{smallmatrix}\right)$ as the product of $v$ applied to the leftmost cell of $\begin{smallmatrix}\beta_{n+2}\left(2i\right)\\
\beta_{n+2}\left(2j\right)
\end{smallmatrix}$ and $v\left(\ell_{i,j}\right)$, where $\ell_{i,j}$ is the binary
lattice formed by deleting the leftmost column of $\begin{smallmatrix}\beta_{n+2}\left(2i\right)\\
\beta_{n+2}\left(2j\right)
\end{smallmatrix}$.

Consider $\beta_{n+2}\left(2k\right)$ for $k=0$, $1$, \ldots ,
$2^{n+1}-1$. The first half of these $2^{n+1}$ strings begin with
$0$ and the second half begin with $1$. Splitting each half in half,
the strings in each of these four parts begin with $0\ 0$, $0\ 1$,
$1\ 0$, and $1\ 1$, in that order. Therefore, as we go down any
column of $Y_{n+1}$ from the top row to the bottom row in four equal
parts, the leftmost cell of $\begin{smallmatrix}\beta_{n+2}\left(2i\right)\\
\beta_{n+2}\left(2j\right)
\end{smallmatrix}$ has upper row $0\ 0$, $0\ 1$, $1\ 0$, then $1\ 1$. Similarly,
going to the right across any row of $Y_{n+1}$, the leftmost cell
has lower row $0\ 0$, $0\ 1$, $1\ 0$, then $1\ 1$. Comparing this
to $X_{n+1}$ as above, the cells evaluated by $v$ are the same in
both $X_{n+1}$ and $Y_{n+1}$.

To complete our proof, we must show for all $i$ and $j$ that $v\left(\ell_{i,j}\right)$
is equal to the entry of $A_{n}$, $B_{n}$, $C_{n}$, or $D_{n}$
that appears in the appropriate block of $X_{n+1}$ as shown above.
Returning to $\beta_{n+2}\left(2k\right)$ for $k=0$, $1$, \ldots ,
$2^{n+1}-1$ and deleting the first digit of each string, this sequence
becomes two copies of the sequence of strings given by $\beta_{n+1}\left(2k\right)$
for $k=0$, $1$, \ldots , $2^{n}-1$. Using this, consider the $2^{n+1}\times2^{n+1}$
matrix whose $\left(i,j\right)$ entry is $v\left(\ell_{i,j}\right)$.
As we go down any column, we get two copies of $\beta_{n+1}\left(2i\right)$
for $i=0$, $1$, \ldots , $2^{n}-1$ in the upper row of $\ell_{i,j}$,
and as we move right across any row, we get two copies of $\beta_{n+1}\left(2j\right)$
for $j=0$, $1$, \ldots , $2^{n}-1$ in the lower row of $\ell_{i,j}$.
Therefore, using the definition of $Y_{n}$, the matrix of $v\left(\ell_{i,j}\right)$
is exactly $\begin{pmatrix}Y_{n} & Y_{n}\\
Y_{n} & Y_{n}
\end{pmatrix}$. By our induction hypothesis, $Y_{n}=X_{n}=\begin{pmatrix}A_{n} & B_{n}\\
C_{n} & D_{n}
\end{pmatrix}$, and thus in the $\left(i,j\right)$ entry of $X_{n+1}$ as given
above, the entry of $A_{n}$, $B_{n}$, $C_{n}$, or $D_{n}$ equals
$v\left(\ell_{i,j}\right)$, as needed.
\end{proof}

\begin{rem*}
For any positive integer $n$ and any nonnegative integers $i$, $j<2^{n}$,
Proposition~\ref{prop: one by n matrix} expresses $v\left(\begin{smallmatrix}\beta_{n+1}\left(2i\right)\\
\beta_{n+1}\left(2j\right)
\end{smallmatrix}\right)$ as an entry of $X_{n}=\begin{pmatrix}A_{n} & B_{n}\\
C_{n} & D_{n}
\end{pmatrix}$. That is, some entry of $A_{n}$, $B_{n}$, $C_{n}$ or $D_{n}$
gives us the evaluation at $v$ of every $\left(2,n+1\right)$ binary
lattice provided that its rightmost column is $\begin{smallmatrix}0\\
0
\end{smallmatrix}$. Although we do not need it, we can instead find $v\left(\begin{smallmatrix}\beta_{n}\left(i\right)\\
\beta_{n}\left(j\right)
\end{smallmatrix}\right)$, i.e., $v$ evaluated at any $\left(2,n\right)$ binary lattice.
For this, we would use the equations of Theorem~\ref{thm: m by n matrix}
with $\hat{A}$ replacing $A$, $\hat{B}$ replacing $B$, etc., except
defining $\hat{A}_{1}=\hat{B}_{1}=\hat{C}_{1}=\hat{D}_{1}=\left(1\right)$.
Adapting our work above, we would then show using induction that $\hat{X_{n}}=\begin{pmatrix}\hat{A}_{n} & \hat{B}_{n}\\
\hat{C}_{n} & \hat{D}_{n}
\end{pmatrix}$ is the $2^{n}\times2^{n}$ matrix whose $\left(i,j\right)$ entry
is $v\left(\begin{smallmatrix}\beta_{n}\left(i\right)\\
\beta_{n}\left(j\right)
\end{smallmatrix}\right)$ for any integer $n\ge2$ and any nonnegative integers $i$, $j<2^{n}$.
Furthermore, by adapting the claim below in the proof of Theorem~\ref{thm: m by n matrix},
the $\left(i,j\right)$ entry of $\hat{X}^{m}_{n}$ equals $\sum_{\ell}v\left(\ell\right)$,
where the sum is over all $\left(m,n\right)$ binary lattices whose
top and bottom rows are equal to $\beta_{n}\left(i\right)$ and $\beta_{n}\left(j\right)$,
respectively. In particular, adding all the entries of $\hat{X}^{m}_{n}$
gives $\sum_{\ell}v\left(\ell\right)$, where the sum is over all
$\left(m,n\right)$ binary lattices.
\end{rem*}
\begin{proof}[Proof of Theorem~\ref{thm: m by n matrix}]
 Fix $n$. We first introduce some terminology. For all positive
integers $m$, we say a \emph{height~$m$ strip} is an $\left(m,n\right)$
binary lattice in which all vertices in the farthest left and farthest
right columns are labeled zero. Then for a height~$m$ strip $\ell$
and a nonnegative integer $i<2^{n-1}$, we say a particular row of
$\ell$ is \emph{determined by $i$} if that row is equal to $0\ \beta_{n-1}\left(i\right)\ 0$.
Note that an $\left(m,n\right)$ framed lattice is exactly a height~$m$
strip in which both the top and bottom rows are determined by $0$.
Thus, to prove our theorem it suffices to show the following.
\begin{claim*}
For all positive integers $m$ and nonnegative integers $i,j<2^{n-1}$,
the $\left(i,j\right)$ entry of $A^{m}_{n}$ equals $\sum_{\ell}v\left(\ell\right)$,
where the sum is over all height~$m$ strips $\ell$ whose first
row is determined by $i$ and whose last row is determined by $j$. 
\end{claim*}
To prove the claim, first observe that for each $i$ and $j$, there
exists a unique height~$1$ strip whose first row is determined by
$i$ and whose last row is determined by $j$. By the second part
of Proposition~\ref{prop: one by n matrix}, $v$ applied to this
strip is equal to the $\left(i,j\right)$ entry of $A_{n}$. This
proves our claim when $m=1$. We continue by induction, so for some
$m\ge2$, assume the claim for $m-1$.

Fix nonnegative integers $i,j<2^{n-1}$, and consider some nonnegative
integer $k<2^{n-1}$. As in the $m=1$ case, the second part of Proposition~\ref{prop: one by n matrix}
tells us that the $\left(i,k\right)$ entry of $A_{n}$ equals $v$
applied to the unique height~$1$ strip with first and last rows
determined respectively by $i$ and by $k$. By our induction hypothesis,
the $\left(k,j\right)$ entry of $A^{m-1}_{n}$ equals $\sum_{\ell}v\left(\ell\right)$,
where the sum is over all height~$m-1$ strips $\ell$ whose first
and last rows are determined respectively by $k$ and by $j$. Consider
applying $v$ both to the height~$1$ strip indicated above and to
one of the height~$m-1$ strips $\ell$ in this sum. Because $v$
is a cell product function, the product of these two numbers is equal
to $v$ applied to the height~$m$ strip formed by overlapping the
bottom row of the height~$1$ strip to the top row of $\ell$, which
we illustrate as

% weird hack for vdot 
\[
v\left(\begin{smallmatrix}0 & \beta_{n-1}(i) & 0\\
0 & \beta_{n-1}(k) & 0
\end{smallmatrix}\right)v\left(\begin{smallmatrix}0 & \beta_{n-1}(k) & 0\\
\vphantom{\int\limits^{x}}\smash{\vdots} & \vphantom{\int\limits^{x}}\smash{\vdots} & \vphantom{\int\limits^{x}}\smash{\vdots}\\
0 & \beta_{n-1}(j) & 0
\end{smallmatrix}\right)=v\left(\begin{smallmatrix}0 & \beta_{n-1}(i) & 0\\
0 & \beta_{n-1}(k) & 0\\
\vphantom{\int\limits^{x}}\smash{\vdots} & \vphantom{\int\limits^{x}}\smash{\vdots} & \vphantom{\int\limits^{x}}\smash{\vdots}\\
0 & \beta_{n-1}(j) & 0
\end{smallmatrix}\right).
\]

\noindent From this we see that the product of the $\left(i,k\right)$
entry of $A_{n}$ and the $\left(k,j\right)$ entry of $A^{m-1}_{n}$
is equal to $\sum_{\ell}v\left(\ell\right)$, where the sum is now
over all height~$m$ strips $\ell$ whose first, second, and last
rows are determined respectively by $i$, by $k$, and by $j$. This
tells us, in turn, that when we additionally sum over all nonnegative
$k<2^{n-1}$, we can conclude that the $\left(i,j\right)$ entry of
$A_{n}\cdot A^{m-1}_{n}=A^{m}_{n}$ is, as desired, the sum $\sum_{\ell}v\left(\ell\right)$
where the sum is over all height~$m$ strips $\ell$ whose first
and last rows are determined respectively by $i$ and by $j$.
\end{proof}


\end{document}
